% Anpassung an Landessprache ---------------------------------------------------
\usepackage[ngerman]{babel}

% Umlaute ----------------------------------------------------------------------
%   Umlaute/Sonderzeichen wie äüöß direkt im Quelltext verwenden (CodePage).
% ------------------------------------------------------------------------------
\usepackage[T1]{fontenc} % Erlaubt automatische Trennung von Worten mit Umlauten.
\usepackage{textcomp} % Euro-Zeichen etc.

% Schrift ----------------------------------------------------------------------
\usepackage{relsize} % Schriftgröße relativ festlegen
\usepackage{lmodern} % bessere Font fuer PDFs als latex default
%\usepackage{libertine} \usepackage{libertinust1math} % eigener fav 1 - text
%\usepackage{mathptmx} % uni vorgabe - times new romes
%\usepackage{nunito} % uni coporate design
%\usepackage[sfdefault]{biolinum}  % eigener fav 1 - überschriften
%% Option 'sfdefault' only if the base font of the document is to be sans serif

% Tabellen ---------------------------------------------------------------------
\PassOptionsToPackage{table}{xcolor}
\usepackage{tabularx}
\usepackage{multicol}
\usepackage{multirow}
\usepackage{bigstrut}
\usepackage{longtable} % für lange Tabellen
\usepackage{array}
\usepackage{ragged2e}
\usepackage{lscape}
\usepackage{makecell} % cellen formatierung mit \makecell
\usepackage{tablefootnote} % erlaubt fotenotes in tabellen ohne tricks
\newcolumntype{w}[1]{>{\raggedleft\hspace{0pt}}p{#1}} % Spaltendefinition rechtsbündig mit definierter Breite

% Grafiken ---------------------------------------------------------------------
\usepackage[dvips,final]{graphicx} % Einbinden von JPG-Grafiken ermöglichen
\usepackage{graphics} % keepaspectratio
\usepackage{floatflt} % zum Umfließen von Bildern
\usepackage{tikz} % visualisation
\usetikzlibrary{trees} % directory trees
\usepackage{wrapfig}

% Sonstiges -------------------------------------------------------------------
\usepackage{amssymb} % fuer /checkmark
\usepackage{amsmath,amsfonts} % Befehle aus AMSTeX für mathematische Symbole
\usepackage{enumitem} % anpassbare Enumerates/Itemizes
\usepackage{xspace} % Leerzeichen hinter parameterlosen Commands einfügen
\usepackage{xparse} % fuer \NewDocumentCommand, ermoeglicht mehrere optionale args
\usepackage{makeidx} % für Index-Ausgabe mit \printindex
\usepackage[useregional]{datetime2} % darstellung eines datum

% Einfache Definition der Zeilenabstände und Seitenränder etc.
\usepackage{setspace}
\usepackage{geometry}
% TODO may include
%\usepackage[compact]{titlesec} % Zeilenabstände
\usepackage{layout} % für command \layout

% Abkürzungsverzeichnis -----------------------------------------------------------
\usepackage[intoc]{nomencl}
\let\abbrev\nomenclature
\renewcommand{\nomname}{Abkürzungsverzeichnis}
\setlength{\nomlabelwidth}{.25\hsize}
\renewcommand{\nomlabel}[1]{#1 \dotfill}
\setlength{\nomitemsep}{-\parsep}

% verzeichnisse -----------------------------------------------------------
\usepackage{url} % URL verlinken, lange URLs umbrechen etc.
\usepackage{chngcntr} % fortlaufendes Durchnummerieren der Fußnoten
\usepackage{ifthen} % bei der Definition eigener Befehle benötigt
\usepackage{todonotes} % definiert u.a. die Befehle \todo und \listoftodos
\usepackage[nolist]{acronym} % liste alle Abkürzungen auf / nur benutzte: [printonlyused]

% literaturverzeichnis --------------------------------------------------------
\usepackage[
    backend=biber, %is default
    style=numeric, % numeric, alphabetic, authoryear, authortitle, science, ieee
    sorting=nyt
]{biblatex}
% nty   Sort by name, title, year.
% nyt   Sort by name, year, title.
% nyvt  Sort by name, year, volume, title.
% anyt  Sort by alphabetic label, name, year, title.
% anyvt Sort by alphabetic label, name, year, volume, title.
% ynt   Sort by year, name, title.
% ydnt  Sort by year (descending), name, title.
% none  Do not sort at all. All entries are processed in citation order.


% PDF-Optionen ----------------------------------------------------------------
\usepackage{pdfpages}
\pdfminorversion=5 % erlaubt das Einfügen von pdf-Dateien bis Version 1.7, ohne eine Fehlermeldung zu werfen (keine Garantie für fehlerfreies Einbetten!)
\usepackage[
    bookmarks,
    bookmarksnumbered,
    bookmarksopen=true,
    bookmarksopenlevel=1,
    colorlinks=true,
% diese Farbdefinitionen zeichnen Links im PDF farblich aus
    linkcolor=main, % einfache interne Verknüpfungen
    anchorcolor=main,% Ankertext
    citecolor=main, % Verweise auf Literaturverzeichniseinträge im Text
    filecolor=main, % Verknüpfungen, die lokale Dateien öffnen
    menucolor=main, % Acrobat-Menüpunkte
    urlcolor=main,
% diese Farbdefinitionen sollten für den Druck verwendet werden (alles schwarz)
%linkcolor=black, % einfache interne Verknüpfungen
%anchorcolor=black, % Ankertext
%citecolor=black, % Verweise auf Literaturverzeichniseinträge im Text
%filecolor=black, % Verknüpfungen, die lokale Dateien öffnen
%menucolor=black, % Acrobat-Menüpunkte
%urlcolor=black,
%
%backref, % Quellen werden zurück auf ihre Zitate verlinkt
    pdftex,
    plainpages=false, % zur korrekten Erstellung der Bookmarks
    pdfpagelabels=true, % zur korrekten Erstellung der Bookmarks
    hypertexnames=false, % zur korrekten Erstellung der Bookmarks
    linktoc=all%linktocpage % Seitenzahlen anstatt Text im Inhaltsverzeichnis verlinken
]{hyperref}


% zum Einbinden von Programmcode ----------------------------------------------
\usepackage{listings}
\usepackage{xcolor}
